%!TEX root = ../../main.tex
% ============================================================
% 几何作图 / 图形放大与缩小
% 本文件由 main.tex 通过 \input 引入，请勿单独编译
% 重构日期: 2026-04-11
% ============================================================

\subsection{解答：按比例放大或缩小图形}

\vspace{0.4cm}

\noindent\textbf{3. (1) 按 1：3 的比画出正方形缩小后的图形。}\\
\noindent\textbf{   (2) 按 2：1 的比画出三角形放大后的图形。}

\vspace{1.2cm}

\begin{center}
% =====================================================================
% 【整体绘制思路：比例缩放作图】
% 1. 基础网格：根据长宽比预估，创建一个宽31, 高10的网格空间。
% 2. 原始黑色图形：
%    - 左侧大正方形：边长 6 格。设定坐标 (2,2) 到 (8,8)。
%    - 右侧小直角三角形：竖直高 3 格，水平宽 2 格。设定坐标 (16,5)-(16,8)-(18,8)。
% 3. 目标红色图形 (red!70!black, line width=1.2pt)：
%    - 正方形按 1:3 缩小：新边长为 6/3 = 2。紧贴在大正方形右侧，顶端对齐。坐标 (10,6) 到 (12,8)。
%    - 三角形按 2:1 放大：新高 3*2 = 6，新宽 2*2 = 4。紧贴在小三角形右侧，顶端对齐。坐标 (20,2)-(20,8)-(24,8)。
% =====================================================================
\begin{tikzpicture}[>=Stealth, scale=0.45]
  % 1. 绘制网格线
  \draw[help lines, dashed, step=1] (0, 0) grid (31, 10);
  \draw[thick, gray!80!black] (0, 0) rectangle (31, 10);
  
  % 2. 绘制原图 (黑色粗线)
  % 大正方形 (6x6)
  \draw[thick, black] (2, 2) rectangle (8, 8);
  % 小直角三角形 (高3, 宽2)
  \draw[thick, black] (16, 5) -- (16, 8) -- (18, 8) -- cycle;
  
  % 3. 绘制按比例缩放的结果图 (深红色加粗)
  % (1) 缩小后的正方形 (1:3 比例 -> 边长 2x2)
  \draw[thick, red!70!black, line width=1.2pt] (10, 6) rectangle (12, 8);
  
  % (2) 放大后的三角形 (2:1 比例 -> 宽4, 高6)
  \draw[thick, red!70!black, line width=1.2pt] (20, 2) -- (20, 8) -- (24, 8) -- cycle;

\end{tikzpicture}
\end{center}

\vspace{0.8cm}

{\small\color{gray}
\textbf{解析：}\\
1. \textbf{图形的缩小}：按 $1:3$ 的比缩小图形，就是把原图形的各边长缩小到原来的 $\frac{1}{3}$。原正方形的边长是 6 格，缩小后得到的新正方形边长应为 $6 \div 3 = 2$ 格。\\
2. \textbf{图形的放大}：按 $2:1$ 的比放大图形，就是把原图形的各边长放大到原来的 2 倍。原直角三角形的竖直直角边长为 3 格，水平直角边长为 2 格。放大后，新的竖直直角边长应为 $3 \times 2 = 6$ 格，新的水平直角边长应为 $2 \times 2 = 4$ 格。最后连接两端点画出新的斜边即可。
}

\vspace{0.6cm}

{\small\color{blue!70!black}
\textbf{绘图基本原理与步骤：}\\
\textbf{1. 统一顶端水平基准线}：经过对参考解答配图的解构，无论是缩小前后的两个正方形，还是放大前后的两个三角形，它们的最高处全部对齐在同一条水平网格线上。因此代码中强制让所有几何图形的高位顶点横向标高均为 \texttt{y=8}（距底边 \texttt{y=0} 共计 \texttt{8} 个单位），以此锚定全图排版的视觉平衡点。\\
\textbf{2. 尺寸方程严格校验}：
   \begin{itemize}
       \item 缩小矩阵：原框 $(2,2) \to (8,8)$，边长跨度为 $6$ 单元。实施算子拉伸 $\times \frac{1}{3}$，生成答案边长跨度为 $2$ 单元。采用 \texttt{rectangle} 指令从 $(10,6)$ 平行绘制至 $(12,8)$。
       \item 放大矩阵：原三角形直角边长为基向量 $(dx=2, dy=3)$。实施放大因子 $\times 2$，其结果图的直角边长强制满足 $(dx=4, dy=6)$。使用全闭合坐标串 \texttt{(20,8) -- (20,2) -- (24,8) -- cycle} 精确构建镜像直角边。
   \end{itemize}
\textbf{3. 专业答卷视觉风格}：所有原题干已知图形统一剥离为黑色常规加粗线型（\texttt{[thick, black]}），所有要求手工绘出的作答结果图层统一覆写辅导读物经典的醒目厚重深红色（\texttt{[thick, red!70!black, line width=1.2pt]}），主次逻辑瞬间明确。
}

\vspace{0.8cm}
\vspace{0.8cm}

\newpage

\subsection{解答：图形的放大与缩小}

\vspace{0.4cm}

\noindent\textbf{5. 按 2:1 的比画出长方形放大后的图形，再按 1:3 的比画出三角形缩小后的图形。}

\vspace{0.8cm}

% =====================================================================
% 【整体绘制思路：图形的放大与缩小】
% 1. 坐标系与网格环境：
%    使用 step=1.0 的虚线在 20x9 的区间绘制方格纸背景，左下角定为坐标原点 (0,0)。
% 2. 比例计算体系：
%    - 原长方形：长 3、宽 1。放大比 2:1 -> 新图长 6、宽 2。
%    - 原三角形：底 9、高 6。缩小比 1:3 -> 新图底 3、高 2。
% 3. 定位布局（复刻参考答案视觉）：
%    - 原长方形置于左上方，新长方形（红框）置于左下方。
%    - 原三角形置于中右方，新三角形（红框）置于其右上方空白处。
% 4. 色彩映射逻辑：
%    - 原题干上的已有图形使用 thick, black 绘制。
%    - 学生需要作答的线条使用 thick, red 绘制，凸显“作图题答案”。
% =====================================================================

\begin{center}
\begin{tikzpicture}[>=Stealth, line width=1.0pt, scale=0.8]

  % =========== 背景网格 ===========
  % \draw[step=1.0, ...]：每隔 1 个单位画一条线
  % grid (20,9)：从 (0,0) 画到 (20,9)，横向 20 个格子，纵向 9 个格子
  % gray!70, dashed：灰色虚线，模拟常见的方格纸效果
  \draw[step=1.0, gray!70, dashed, very thin] (0,0) grid (20,9);
  % 画网格最外侧的实线边框
  \draw[gray!80, thin] (0,0) rectangle (20,9);

  % =========== 原图形（题干黑色图形） ===========
  % 1. 原始长方形 (长 3, 宽 1)
  % 放在左上角，坐标从 (1,7) 到 (4,8)
  \draw[black] (1,7) rectangle (4,8);

  % 2. 原始直角三角形 (底 9, 高 6)
  % 左侧直角边对齐 x=9，顶部对齐 y=8 (与长方形顶端齐平)
  \coordinate (T_top) at (9,8);     % 顶点，位于上方
  \coordinate (T_bl)  at (9,2);     % 直角顶点，位于下方
  \coordinate (T_br)  at (18,2);    % 锐角顶点，位于右侧 (9+9=18)
  % cycle 自动连接 (18,2) 并闭合至起点 (9,8)
  \draw[black] (T_top) -- (T_bl) -- (T_br) -- cycle;


  % =========== 作答图形（红色缩放图形） ===========
  % 3. 放大后的长方形：按 2:1 放大 (长 6, 宽 2)
  % 放置在原长方形下方（间隔 1 格），坐标从 (1,4) 到 (7,6)
  \draw[red] (1,4) rectangle (7,6);

  % 4. 缩小后的三角形：按 1:3 缩小 (底 3, 高 2)
  % 放置在原三角形的右上方，顶部平起，即顶部 y=8
  \coordinate (T_ans_top) at (14,8);    % 顶点，y=8
  \coordinate (T_ans_bl)  at (14,6);    % 直角顶点，高 2，所以 y=8-2=6
  \coordinate (T_ans_br)  at (17,6);    % 锐角顶点，底 3，所以右侧 = 14+3=17
  \draw[red] (T_ans_top) -- (T_ans_bl) -- (T_ans_br) -- cycle;

\end{tikzpicture}
\end{center}

\vspace{0.6cm}

{\small\color{gray}
\textbf{解析：}\\
1. **长方形的放大**：测量原长方形可知其长为 $3$ 个方格，宽为 $1$ 个方格。按照 $2:1$ 的比放大，也就是将边长扩大到原来的 $2$ 倍。放大后的长为 $3 \times 2 = 6$ 个方格，宽为 $1 \times 2 = 2$ 个方格。据此在大网格中画出新的 $6 \times 2$ 长方形。\\
2. **三角形的缩小**：测量原直角三角形可知其两条直角边（底和高）分别为 $9$ 个方格和 $6$ 个方格。按照 $1:3$ 的比缩小，即边长缩小为原来的 $\frac{1}{3}$。缩小后的底为 $9 \div 3 = 3$ 个方格，高为 $6 \div 3 = 2$ 个方格。选取网格右侧空白处，以 $3$、 $2$ 为直角边画出一个符合原图姿态的直角三角形。
}

\vspace{0.6cm}

{\small\color{blue!70!black}
\textbf{绘图基本原理与步骤：}\\
\textbf{1. 虚线网格系统搭建}：运用 \texttt{grid} 命令配合 \texttt{dashed} 属性，构建底层的 $20 \times 9$ 网格，为线段测量和比例计算提供像素级准确的参照系。\\
\textbf{2. 解构题设与变量转换}：首先固定题干已给图形（\texttt{black} 实线），获取原尺寸（$3\times 1$ 与 $9\times 6$）；随后分别进行代数倍乘与除法（放大倍数与缩小比例），得出目标尺寸（$6\times 2$ 与 $3\times 2$）。\\
\textbf{3. 多元素布局防重叠}：通过对全局画布（$0 \le x \le 20, 0 \le y \le 9$）进行预规划：将大面积答案（大红框）沉于左侧中部，小面积答案（小红三角）悬于右上，既满足紧凑要求亦实现与参考答案在视觉落点上的高度一致。\\
\textbf{4. 定向颜色标注}：严格施行“原题黑、答案红”的双色区分法，强化辅助教材的直观导向性。
}

\newpage
